\documentclass[11pt]{ctexart}
\usepackage{amsmath,amssymb,enumitem,geometry}
\geometry{margin=1in}
\setlist[enumerate]{leftmargin=*}
\title{Week 2 Solutions: Algorithmic Thinking, Data Structures, and Proof Patterns}
\author{TCS Self-Study OS}
\date{Generated 2026-05-06}

\begin{document}
\maketitle

每题包含 proof strategy 或 solution idea，并列出 common mistake。算法题分 correctness 和 complexity；randomized 题说明 sample space；amortized 题说明方法。

\section{Day 1 Solutions}
\begin{enumerate}
  \item \textbf{Problem 1.} Strategy: 先定义 input domain 和 output relation。Solution: Input 是长度为 n 的可比较元素序列 A。Output 是序列 B，满足 B 是 A 的 permutation，且对所有 1<=i<n 有 B\_i <= B\_{i+1}。Precondition 是元素之间有 total preorder 或 total order；postcondition 是 permutation 加 sortedness。n 是元素数。Common mistake: 只写“输出排好序数组”，没有 permutation 条件。
  \item \textbf{Problem 2.} Strategy: 用常见算法区分两个性质。Solution: 标准 merge sort 可 stable，但需要 O(n) auxiliary space，因此不是严格 O(1) in-place。标准 heapsort 通常 in-place，但相等 key 的记录可能交换相对顺序，因此不 stable。Common mistake: 以为 stable 描述空间，in-place 描述相等元素。
  \item \textbf{Problem 3.} Strategy: initialization, maintenance, termination。Solution: i=1 时单元素前缀满足 sortedness 和 permutation。维护时，while 只把大于 key 的元素右移，不改变前缀多重集合；找到位置后插入 key，左边 <= key，右边 > key，因此 A[0..i+1) sorted。终止 i=n 得整个数组 sorted permutation。Common mistake: 只说“插入到正确位置”但不证明 permutation 保持。
  \item \textbf{Problem 4.} Strategy: 每层成本求和。Solution: 第 l 层有 2\textasciicircum{}l 个子问题，每个大小 n/2\textasciicircum{}l，merge 总成本为 2\textasciicircum{}l c(n/2\textasciicircum{}l)=cn。深度 log\_2 n。叶子总成本为 dn。总成本 cn log\_2 n + dn = O(n log n)。Common mistake: 忽略叶子成本或默认 n 必为 2 的幂；一般 n 可用上下取整处理。
  \item \textbf{Problem 5.} Strategy: 不同排列需要不同输出。Solution: 对两两不同元素，存在 n! 种可能相对顺序。若两个不同排列到达同一个 leaf，deterministic algorithm 会输出同一个 permutation，但正确排序要求两个输入的相对顺序不同，所以至少一个输出错误。因此每个排列需要独立 leaf，leaves >= n!。Common mistake: 把输入值数量和相对顺序数量混淆。
  \item \textbf{Problem 6.} Strategy: 检查 model。Solution: comparison lower bound 只适用于只能通过比较获得顺序信息的算法。Counting sort 使用 key 的整数范围和数组计数，利用了非比较操作和 key universe 假设，因此不在 comparison model 内。Common mistake: 把所有 sorting algorithms 都放进 comparison model。
  \item \textbf{Problem 7.} Strategy: 合并同一模型中的 upper/lower bound。Solution: 存在 comparison sorting algorithm merge sort 达到 O(n log n)，且任何 comparison sorting algorithm worst-case 至少 Omega(n log n)，所以 comparison sorting problem 的 optimal worst-case comparison complexity 是 Theta(n log n)。Common mistake: 对单个算法说 lower bound，而不是对问题/模型说 optimal complexity。
\end{enumerate}

\section{Day 2 Solutions}
\begin{enumerate}
  \item \textbf{Problem 8.} Strategy: specification first。Solution: Precondition: A sorted, lo=0, hi=n-1。Postcondition: 返回某个 index i with A[i]=x，或证明不存在。Invariant: 若 x 在原数组中，则至少一个出现位置仍在当前区间 [lo,hi]。Common mistake: 忘记 sortedness precondition。
  \item \textbf{Problem 9.} Strategy: 展开或 recursion tree。Solution: n 是 input size。每层成本 1，规模折半到 1 需要 log\_2 n 层，所以 T(n)=Theta(log n)。Common mistake: 写 O(log) 而不写变量。
  \item \textbf{Problem 10.} Strategy: induction on n。Solution: n<=1 成立。假设所有小于 n 的数组都被正确排序。算法递归排序左右两半，由 IH 得 L,R 是各自 sorted permutations。merge lemma 保证 merge(L,R) 是 L union R 的 sorted permutation，也就是 A 的 sorted permutation。Common mistake: 证明了 sortedness 但没有证明 permutation。
  \item \textbf{Problem 11.} Strategy: 猜 T(n)<=c n log\_2 n。Solution: 对 n>=2，T(n)<=2 c(n/2)log(n/2)+n = c n log n - c n + n <= c n log n，当 c>=1。base 可调大 c 覆盖。Common mistake: induction 假设没有写适用范围。
  \item \textbf{Problem 12.} Strategy: 比较 output specification。Solution: Selection 只要求第 k 小元素，不要求输出所有元素的 total order。Baseline 是先 sort 再取第 k 个，时间 O(n log n)，但 specification 允许更快算法如 expected O(n) quickselect。Common mistake: 把排序作为 selection 的一部分。
  \item \textbf{Problem 13.} Strategy: 定义 middle-half pivot。Solution: 若 pivot rank 在 [n/4,3n/4]，递归子问题规模至多 3n/4，称为 good。随机 uniform pivot 落在这一区间的概率至少 1/2。因此期望经过常数次 partition 得到一次规模显著下降，累计线性层级成本为 O(n)。Common mistake: 把这个直觉当完整 tail proof。
  \item \textbf{Problem 14.} Strategy: 子问题返回 pair。Solution: 分成两半递归求 (min,max)，combine 用两次比较合并。T(n)=2T(n/2)+O(1)，所以 T(n)=O(n)。Correctness by induction: 子 pair 对子数组正确，combine 取两个 min 的较小者和两个 max 的较大者。Common mistake: combine step 没有 specification。
\end{enumerate}

\section{Day 3 Solutions}
\begin{enumerate}
  \item \textbf{Problem 15.} Strategy: interface 不提 heap array。Solution: insert(Q,x) 返回包含旧元素加 x 的 priority queue；extract-min(Q) 在 Q 非空时返回最小 key 元素并从 Q 删除一个该元素。Precondition for extract-min: Q nonempty。Common mistake: 在 interface 中写 parent/child array details。
  \item \textbf{Problem 16.} Strategy: 检查每个 parent-child edge。Solution: parent 1 <= 3,2；3 <= 7,5；2 <= 4，所以满足 min-heap invariant。Common mistake: 要求整个左子树小于右子树。
  \item \textbf{Problem 17.} Strategy: violation path。Solution: append 新元素后，只有新叶到 root 的路径可能违反 parent<=child。每次与 parent swap 后，当前 child 位置满足与其 children 的关系，因为较大 parent 被放下且不小于原 child 的 children；路径外边不变。循环停止时 root 或 parent<=child，全树 invariant 恢复。Common mistake: 不说明路径外边为什么不受影响。
  \item \textbf{Problem 18.} Strategy: bucket length expectation。Solution: 查询 key 的 bucket 中包含该 key 和其他碰撞 keys。每个其他 key 进入同一 bucket 的概率 1/m，期望碰撞数 (n-1)/m <= alpha。遍历成本为 1 加期望碰撞数，所以 O(1+alpha)。Common mistake: 忘记这是 expectation under hashing assumption。
  \item \textbf{Problem 19.} Strategy: resize sizes are powers of two。Solution: 扩容时复制 1,2,4,...,2\textasciicircum{}r 个元素，最大复制规模小于 m。几何和 <2m。普通写入 m 次，总成本 O(m)，amortized O(1)。Common mistake: 只看某次 resize O(n) 就否定 amortized O(1)。
  \item \textbf{Problem 20.} Strategy: 选 Phi=2n-c after at least half-full invariant。Solution: 在 load 维持足够高的模型中可用 Phi=2n-c。普通 append 增加实际成本 1、potential 增加 2；积累信用。满时 n=c，resize 到 2c 后 append，复制 c 个元素，但 potential 从约 c 降到常数级，释放 O(c) 信用支付复制。Common mistake: potential 取负而未修正。
  \item \textbf{Problem 21.} Strategy: interface/spec/assumption。Solution: Interface: insert, find, delete by key。Chaining representation invariant: key k 存在于 bucket h(k)。Worst-case 所有 keys collision，operation O(n)。Under uniform hashing and bounded alpha，expected O(1)。Common mistake: 对 adversarial keys 仍声称 worst-case O(1)。
\end{enumerate}

\section{Day 4 Solutions}
\begin{enumerate}
  \item \textbf{Problem 22.} Strategy: count storage。Solution: adjacency list 用 O(|V|+|E|) 存 headers 和 edges；adjacency matrix 用 O(|V|\textasciicircum{}2) entries。Dense graph 中 matrix 可接受，sparse graph 中 list 更省。Common mistake: 对所有图都写 O(|E|)。
  \item \textbf{Problem 23.} Strategy: layer expansion。Solution: dist(1)=0, dist(2)=1, dist(3)=2, dist(4)=3；parents: parent(2)=1, parent(3)=2, parent(4)=3。Common mistake: parent tree 与原图所有边混淆。
  \item \textbf{Problem 24.} Strategy: layer invariant。Solution: BFS queue 按非降 distance 处理。若 v 第一次由 u 发现，则存在 path length dist[u]+1。若有更短 path，其前驱 w distance 更小，会在 u 之前或同层更早处理并发现 v，矛盾。因此 dist[v] shortest。Common mistake: 未说明 queue order。
  \item \textbf{Problem 25.} Strategy: construct long branch first。Solution: 图有 edges s-a, a-b, b-t, s-t。若 DFS 从 s 先访问 a，则找到 s-a-b-t 长度 3；shortest path 是 s-t 长度 1。Common mistake: 把 DFS reachability 与 shortest path 混淆。
  \item \textbf{Problem 26.} Strategy: outer loop over vertices。Solution: 对每个 unvisited vertex 启动 BFS/DFS，访问到的集合是一个 component。每个 vertex 入队/入栈一次，每条 edge 被检查常数次，复杂度 O(|V|+|E|)。Common mistake: 只从一个 source 出发。
  \item \textbf{Problem 27.} Strategy: edge finish relation。Solution: 对任意 edge u->v。若 DFS 从 u 发现 v，则 v 完成后 u 才完成，finish(u)>finish(v)。若 v 已完成，则 finish(v)<finish(u)。若 v 在递归栈中则有 cycle，与 DAG 矛盾。因此 u 总在 v 前。Common mistake: 不排除 back edge。
  \item \textbf{Problem 28.} Strategy: settled invariant breaks。Solution: Dijkstra 的 invariant 是每次取出的最小 tentative distance 已最终确定。若存在负权边，之后可能通过尚未处理 vertex 把已确定 vertex 的距离降低，破坏 invariant。Common mistake: 认为只要没有 negative cycle 就可直接用 Dijkstra。
\end{enumerate}

\section{Day 5 Solutions}
\begin{enumerate}
  \item \textbf{Problem 29.} Strategy: compare greedy vs optimal。Solution: Greedy 选 4，剩 2，得 4+1+1 共 3 枚；最优为 3+3 共 2 枚。Common mistake: 用 canonical coin systems 的经验推广到任意系统。
  \item \textbf{Problem 30.} Strategy: sort by finish。Solution: finish order: [1,2), [0,3), [2,4), [3,5)。选 [1,2)，再选 [2,4)，不能选 [3,5) 因为与 [2,4) 重叠。输出 2 个 intervals。Common mistake: 按 start time 选。
  \item \textbf{Problem 31.} Strategy: replace first optimal interval。Solution: 令 g 为最早结束 interval，O 为任意 optimal solution，o 为 O 中最早结束 interval。因为 f\_g <= f\_o，用 g 替换 o 不会与 O 中其他 intervals 冲突，且大小不变，所以存在包含 g 的 optimal solution。Common mistake: 说所有 optimal 都包含 g。
  \item \textbf{Problem 32.} Strategy: compare kth finish times。Solution: 令 g\_k 是 greedy 第 k 个 interval，o\_k 是任一 optimal solution 的第 k 个 interval。证明 f(g\_k) <= f(o\_k) 对所有 k 成立。若 optimal 有 r 个 intervals，greedy 至少也可安排到 r 个。Common mistake: 比较 start time 而不是 finish time。
  \item \textbf{Problem 33.} Strategy: exchange equal weight。Solution: 若某解取了低 density item 的正重量，同时高 density item 未取满，则把一小段重量从低 density 换到高 density，价值不降且若 density 严格更高则增加。重复得到按 density 的 greedy form。Common mistake: 套用到 0/1 knapsack。
  \item \textbf{Problem 34.} Strategy: component cut。Solution: 当前 forest 的某个 component S 与 V\S 形成 cut。Kruskal 选择的连接该 component 到外部的最轻安全边跨越此 cut。根据 cut property，该 light edge 属于某个 MST，可加入而不破坏存在最优扩展。Common mistake: 未说明 cut 如何由 component 给出。
  \item \textbf{Problem 35.} Strategy: 一个长 interval 堵住多个短 interval。Solution: intervals [0,10), [1,2), [2,3), [3,4)。Earliest start 选 [0,10) 得 1；earliest finish 可选三个短 interval。Common mistake: 只凭直觉说开始早更好。
\end{enumerate}

\section{Day 6 Solutions}
\begin{enumerate}
  \item \textbf{Problem 36.} Strategy: avoid ambiguity。Solution: dp[i] 表示以 A[i] 结尾的 longest increasing subsequence length。最终答案是 max\_i dp[i]。Common mistake: 以为 dp[n-1] 就是答案。
  \item \textbf{Problem 37.} Strategy: empty prefix。Solution: dp[i][0]=i，因为把长度 i 的 string 变空需要 i 次 delete；dp[0][j]=j，因为从空变成长度 j 的 string 需要 j 次 insert。Common mistake: base case 全设为 0。
  \item \textbf{Problem 38.} Strategy: previous element case split。Solution: dp[i]=1+max{dp[j]:j<i,A[j]<A[i]}，若集合空则 dp[i]=1。任何以 i 结尾的 increasing subsequence 的前一个 index 必为 j<i 且 A[j]<A[i]；反之任意这样的 subsequence 可接 A[i]。Common mistake: 用 A[j]<=A[i] 改变问题为 nondecreasing。
  \item \textbf{Problem 39.} Strategy: last operation。Solution: dp[i][j]=min(dp[i-1][j]+1, dp[i][j-1]+1, dp[i-1][j-1]+cost)，其中 cost=0 if x[i-1]=y[j-1] else 1。三项对应 delete x last, insert y last, match/substitute。Common mistake: index i 与字符 x[i-1] 混淆。
  \item \textbf{Problem 40.} Strategy: dependencies smaller。Solution: dp[i][j] 依赖 dp[i-1][j], dp[i][j-1], dp[i-1][j-1]。按行主序填时，上一行和左侧当前行已计算，依赖满足。Common mistake: 使用未计算状态。
  \item \textbf{Problem 41.} Strategy: include/exclude item。Solution: dp[i][w]=dp[i-1][w] if weight\_i>w，否则 max(skip, value\_i+dp[i-1][w-weight\_i])。时间 O(nW)，但 W 的二进制编码长度是 log W，所以该复杂度不一定是输入长度的 polynomial。Common mistake: 忽略数字编码。
  \item \textbf{Problem 42.} Strategy: parent pointers。Solution: 计算 dp[i] 时记录达到最大值的前驱 parent[i]。从 max dp 的 index 反复跟 parent 回溯，再 reverse。Common mistake: 只计算 value 却声称已经输出解。
\end{enumerate}

\section{Day 7 Solutions}
\begin{enumerate}
  \item \textbf{Problem 43.} Strategy: correctness vs time。Solution: Las Vegas 总输出正确，时间随机，如 randomized quickselect 若一直运行到确定 rank。Monte Carlo 时间受控但可错，如抽样 majority toy test。Common mistake: 认为 randomized 就一定可能错。
  \item \textbf{Problem 44.} Strategy: random choices sequence。Solution: sample space 是每次递归 pivot rank 或 pivot index 的所有可能序列，概率由每层 uniform pivot choice 决定。Running time 是该 sample space 上的 random variable。Common mistake: 把输入分布当成算法随机性。
  \item \textbf{Problem 45.} Strategy: quantify randomness vs input。Solution: Expected time 固定输入后对算法 random bits 取期望；worst-case time 对输入和随机选择中的最坏路径取最大。Randomized quickselect expected O(n)，但某些 pivot 序列导致 O(n\textasciicircum{}2)。Common mistake: 认为 expected O(n) 表示所有运行都 O(n)。
  \item \textbf{Problem 46.} Strategy: independence product。Solution: sample space 是 r 次 independent random bits。整体 failure 事件是每次都失败，概率 <=(1/3)\textasciicircum{}r。Common mistake: 没有独立性就相乘。
  \item \textbf{Problem 47.} Strategy: count bad runs。Solution: 整体错误当且仅当超过 r/2 次 base run 错误。可用 binomial tail 或 Chernoff bound 控制。Common mistake: 以为只要有一次正确就 majority 正确。
  \item \textbf{Problem 48.} Strategy: correctness conditional on any pivot。Solution: 对任意 pivot，partition 得 L,E,G 且所有 L<=E<=G。递归排序 L,G 后 concatenation sorted permutation。随机性只影响运行时间分析，不影响 correctness。Common mistake: 把随机 pivot 当成正确性必要条件。
  \item \textbf{Problem 49.} Strategy: define random hash function。Solution: 若 h(k) uniform over m buckets 且 pairwise collision probability 1/m，则 Pr[h(x)=h(y)]=1/m。Bucket length 的 expectation 用 indicators I\_y=1[h(y)=h(x)] 求和。Common mistake: 不说明 hash function 的随机模型。
\end{enumerate}

\section{Self-Test Answer Key}
\begin{enumerate}
  \item Short answers should define the requested term, state the relevant model or assumption, and give one example or non-example.
  \item Correctness proofs are graded for specification, invariant or induction hypothesis, maintenance/step argument, termination, and edge cases.
  \item Complexity analyses must name variables and distinguish worst-case, expected, and amortized claims.
  \item Design questions must include problem statement, input, output, algorithm idea, pseudocode, correctness, and complexity.
  \item Randomized answers must define the sample space, random variables, bad event, and independence assumptions.
  \item The comprehensive problem should combine a clean specification, at least one proof pattern, and a final complexity statement.
\end{enumerate}

\section{Self-Test Solution Outlines}

\subsection*{Short Answer}
\begin{enumerate}
  \item An algorithmic problem is a family of valid instances together with acceptable outputs. An algorithm is a procedure proposed to solve that problem.
  \item Binary search precondition: array is sorted and finite. Postcondition: return an index containing the target, or report that no such index exists.
  \item The comparison model allows pairwise comparisons and forbids using numeric key structure such as counting buckets unless that operation is part of the model.
  \item Stability preserves relative order of equal keys. In-place concerns auxiliary space. Neither property implies the other.
  \item A representation invariant is a condition all valid internal states of a data structure must satisfy.
  \item Amortized analysis bounds a worst-case sequence of operations. Average-case analysis requires an input or randomness distribution.
  \item BFS invariant: vertices are processed by nondecreasing distance layers, and first discovery gives a shortest unweighted path length.
  \item Greedy choice is safe means there exists an optimal solution containing the greedy choice, so committing to it does not reduce optimal value.
  \item For edit distance, dp[i][j] is the minimum number of edits turning the first i characters of x into the first j characters of y.
  \item Las Vegas algorithms are always correct but may have random running time. Monte Carlo algorithms may err with bounded probability.
\end{enumerate}

\subsection*{Algorithm Correctness Proofs}
\begin{enumerate}
  \item Use the sorted-prefix invariant: before outer iteration i, A[0..i) is a sorted permutation of the original prefix. Prove initialization, maintenance by shifting, and termination at i=n.
  \item Induct on n. Recursive calls sort both halves by induction. The merge lemma says merging two sorted permutations gives a sorted permutation of their union.
  \item After appending, only the new leaf-to-root path can violate heap order. Each swap moves the violation upward and preserves off-path edges. Stop condition restores parent-child order everywhere.
  \item Use the BFS layer invariant. If v is first discovered from u, then a path of length dist[u]+1 exists. A shorter path would have a predecessor in an earlier layer, processed before u.
  \item Use exchange argument. Replace the first interval of an optimal solution with the earliest-finishing interval; feasibility and cardinality are preserved, then recurse on the remaining compatible intervals.
\end{enumerate}

\subsection*{Complexity Analysis}
\begin{enumerate}
  \item In T(n)=2T(n/2)+n, each recursion-tree level costs n and there are log n levels, so T(n)=Theta(n log n).
  \item Binary search halves the candidate interval each comparison, so worst-case comparisons are Theta(log n).
  \item Chained hash search is expected O(1+alpha) under simple uniform hashing, where alpha=n/m. Worst case is O(n) if all keys collide.
  \item Dynamic array doubling copies 1+2+4+...+less than m elements over m appends, so total cost is O(m) and amortized cost is O(1) by aggregate method.
\end{enumerate}

\subsection*{Design Questions}
\begin{enumerate}
  \item Min/max divide-and-conquer returns a pair for each half and combines by comparing two minima and two maxima. Correctness is induction on subarray length; time satisfies T(n)=2T(n/2)+O(1)=O(n).
  \item LIS design: dp[i] is LIS length ending at i; parent[i] stores predecessor. Transition checks all j<i with A[j]<A[i]. Reconstruction follows parent from an index with maximum dp value.
  \item Components design: scan all vertices; start BFS or DFS at each unvisited vertex; the reached set is one component. Each vertex and edge is processed a constant number of times, so time is O(|V|+|E|).
\end{enumerate}

\subsection*{Randomized and Amortized Questions}
\begin{enumerate}
  \item Randomized quickselect sample space is the sequence of pivot choices. Running time is the number of comparisons or partition work as a random variable over that sample space.
  \item Repeat the Monte Carlo algorithm r independent times. If overall failure means all runs fail, probability is at most (1/4)\textasciicircum{}r. If majority vote is used, the bad event is more than half the runs failing.
\end{enumerate}

\subsection*{Comprehensive Question}
For shortest paths in a DAG, topologically sort reachable vertices from s, initialize dist[s]=0 and other distances to infinity, then relax outgoing edges in topological order. Correctness is induction over topological order: when vertex u is processed, all possible predecessors of u have already been processed, so dist[u] is final. Each edge is relaxed once after topological ordering, so time is O(|V|+|E|) with adjacency lists, and space is O(|V|+|E|) for the graph plus O(|V|) for distances and parents.

\end{document}
